\section{Isolated information of one stored pair}
\label{sec:h-store}

The open row asks for an entropy built from exterior-blind store entries.
The pair data \((s,d)\) already determine a single real weight.
Lemma~\ref{lem:bond-expand} gives
\[
P_-=\frac12-\frac{B_{ij}}{q_{ij}}
\]
whenever \(q_{ij}>0\), and Corollary~\ref{cor:Bzero} identifies the maximum of
\(h_2(P_-)\) with a norm-neutral bond.
Remark~\ref{rem:h2} keeps \(h_2\) optional: the ledger has not introduced
\(-\sum p\log p\).

Theorem~\ref{thm:haar}, Proposition~\ref{prop:arcsine} and
Proposition~\ref{prop:ground} are three exact means on one pair:
\(1/(2\ln 2)\), \(2-1/\ln 2\), and \(0\).
Proposition~\ref{prop:oddsat} is the same zero on the opposite endpoint.
Every sheet-odd state has \(h_2=0\) on each occupied rung, and the sum of
the rung bonds of a normalised state is \(-1/2\).
Proposition~\ref{prop:lambda} is the scale-separated limit, in which
\(h_2\to 1\) because the amplitudes differ, not because a phase was selected.
The same \(P_-\) is the daughter weight in the split fibre
(Section~\ref{sec:h-split}) and the summand in any cut sum that one might
compare with the I1 bit (Section~\ref{sec:bond}).
The continuous dimension of \(d\) is Section~\ref{sec:kernel} and is not this sum.
On a matching of disjoint pairs the optional bits are independent, with
mean and variance given by Theorem~\ref{thm:disjoint}.
A shared vertex keeps the one-edge mean and produces the positive
covariance of Proposition~\ref{prop:cov}.
Either way \(h_2\) is constant on the circle of
Proposition~\ref{prop:circle}, so the bit does not resolve the current.
Theorem~\ref{thm:condJ} gives the residual law of \(\mathcal{J}/q\) given
\(B/q\), including its differential entropy.
That entropy is continuous and is not a term in \(h_2\).
